% ============================================================
%  文档类选项（在 main.tex 的 \documentclass[...]{tongjithesis} 中设置）
% ============================================================

% ============================================================
%  封面信息
% ============================================================
\school{计算机科学与技术学院}
\major{软件工程}
\student{2251076}{刘嘉程}
\thesistitle{基于Spring Boot的校内综合交通平台设计与实现}{}
\thesistitleeng{Design and Implementation of a Campus Comprehensive Transportation Platform Based on Spring Boot}{}
\thesisadvisor{唐剑锋}
\thesisdate{2026}{5}{1}

% ============================================================
%  信息说明页
% ============================================================
\infotype{design}

\infoabstract{%
本文围绕高校多校区办学背景下师生出行信息分散、拼车组织效率偏低及校内交通服务缺乏统一入口的问题，设计并实现了一个基于Spring Boot的校内综合交通平台。系统采用前后端分离与微服务架构，以微信小程序作为统一入口，整合拼车、论坛、聊天和公共交通信息查询等功能，并结合MySQL与MongoDB实现混合存储。在拼车业务中，系统提出“规则筛选+随机森林排序”的二层匹配机制，以提升候选订单排序效果。测试结果表明，平台能够稳定完成核心业务流程，具有较好的实用性与扩展性。%
}

\infodrawings{0}
\infowordcount{23000}

\infomaterials{
  \item 校内综合交通平台项目源代码
}
